\scanpage{8}
induit des morphismes \inserted{sur} $A'\times_S e$ et $e\times_S A$ qui sont triviaux. Donc il est trivial sur \deleted{\unsure{$S$|$\xi$}} $A\times A'$, \deleted{\unsure{d'après}} étant \unsure{multiplicatif} sur $A\times A'$, qui est un schéma abélien.
% slanted
\margin{ici on utilise le fait que $A'$ est \uline{abélien}}
Donc on est ramené à prouver que le morphisme $(**)$ \deleted{\unsure{est une}} a un noyau nul, et il suffit d'ailleurs de prouver que \deleted{\unsure{il}} c'est une immersion fermée. Comme $\underline{\text{\unsure{NérSév|NérSé}}}_{A/S}$ \deleted{est} donc \deleted{\unsure{tout}} \unsure{son -préschéma|son sous-préschéma} est non ramifié sur $S$, on est ramené à prouver \deleted{que les} pour le premier pt que les fibres du noyau de $\underline{\text{NérSév}}_{A/S}\to\underline{\mathrm{Hom}}_{S\text{-gr}}(A,A')$ sont ensemblistement réduites à un pt. Donc on est ramené au cas où $S$ est le spectre d'un corps alg. clos, et à prouver ceci :

\noindent\llap{\smash{\rule[-3.7ex]{0.4pt}{4.2ex}}\quad}\uline{Lemme} $A$ var. \unsure{abélienne|abélien} sur un corps $k$ \unsure{alg. clos|alg clos} $\xi$ une classe de diviseurs sur $A$, \unsure{invariante|invariant} par translation. Alors $\xi$ est alg. équivalente à $0$

C'est connu, et sera prouvé plus bas \ldots..
% archive sheet number: 7

